%%%PAGE 143 rot=0 legibility=good summary=卫星几何类习题续：A/B两卫星可直接微波通信的时间；卫星绕木星N圈求轨道半径与木星第一宇宙速度(9)；由张角和周期求地球密度
%%%BLOCK id=143-01 ch=05 sec=追及相遇·卫星通讯 kind=problem alt=none cont=none y=0.03-0.31
% 页面右上角另一页露出的字迹“-1.429 g·L^-1”不属于本页，忽略
\textbf{作切线\unclear{求}}

\begin{problem}
$A$、$B$两卫星某时刻相距最近，若该时刻开始至两卫星再次相距最近的时间为$T$，在这段时间$T$内两卫星可直接利用微波通信的时间为% 行末贴近页边，后面可能还有字
%FIG p143_a 0.08 0.09 0.30 0.25
\notefig[0.3]{figures/p143_a.png}
\begin{solution}
\[ \omega_A t-\omega_B t=2\theta\ \text{①}\qquad \omega_A T-\omega_B T=2\pi\ \text{②}\qquad \dfrac{\text{①}}{\text{②}}\text{得}\ \dfrac{t}{T}=\dfrac{67}{180^\circ} \]
\[ \sin\alpha=\dfrac{R}{R+\frac{2}{3}R}=\dfrac{3}{5}\qquad \alpha=37^\circ \]
\[ \sin\beta=\dfrac{R}{R+R}=\dfrac{1}{2}\qquad \beta=30^\circ \]
时间$t$应大于$\dfrac{67}{180}T$　（不能等于）
\end{solution}
\end{problem}
%%%END
%%%BLOCK id=143-02 ch=05 sec=卫星几何·张角问题 kind=problem alt=none cont=none y=0.31-0.635
\begin{problem}[9]
\textbf{9、}卫星$t$ s绕木星$N$圈，速率为$V$，求轨道半径和木星第一宇宙速度
%FIG p143_b 0.04 0.318 0.28 0.505
\notefig[0.3]{figures/p143_b.png}
\begin{solution}
由$t=\dfrac{N\,2\pi r}{V}$　$\therefore\ r=\dfrac{Vt}{2\pi N}$

\unclear{表面}$\dfrac{GMm}{R^2}=\dfrac{mV_1^2}{R}$，　$V_1=\sqrt{\dfrac{GM}{R}}$
\[ \begin{cases}
\text{卫星}\ \dfrac{GMm'}{r^2}=\dfrac{m'V^2}{r}\\[4pt]
R=r\sin\dfrac{\theta}{2}
\end{cases} \]
\[ GM=V_1^2R \]
\[ \therefore\ V_1=\dfrac{V}{\sqrt{\sin\dfrac{\theta}{2}}} \]
\end{solution}
\end{problem}
%%%END
%%%BLOCK id=143-03 ch=05 sec=天体质量密度·张角法 kind=problem alt=none cont=none y=0.635-1.00
\begin{problem}
求地球密度
%FIG p143_c 0.13 0.685 0.31 0.875
\notefig[0.25]{figures/p143_c.png}
\begin{solution}
\[ \rho=\dfrac{3\pi}{GT^2} \]
\[ \left(\dfrac{R_{\text{地}}}{r}\right)^3=\left(\dfrac{T_{\text{地}}}{T}\right)^2 \]
% 以下两项为原稿右侧划去的草算
\struck{$R_{\text{地}}$}　\struck{\unclear{$\left(\dfrac{R_{\text{地}}}{r}\right)^3=\dfrac{T_{\text{地}}^2}{T_{\text{地}}^2}$}}
\[ \dfrac{R_{\text{地}}}{r}=\sin\dfrac{\theta}{2} \]
\[ \therefore\ \left(\sin\dfrac{\theta}{2}\right)^{\frac{3}{2}}T=T_{\text{地}} \] % 原稿 T 处有涂改
\[ \rho_{\text{地}}=\dfrac{3\pi}{G\sin^{3}\dfrac{\theta}{2}\,T^2} \]
\end{solution}
\end{problem}
%%%END
%%%PAGEEND 143
